\subsection{Regulatory Basis and the Form of the Requirement}
\label{sec:interference-basis}

In the United States no numerical emission limit applies to high-voltage
transmission equipment in the PLC and RI bands. Overhead lines and converter
stations are incidental radiators under 47~CFR Part~15, \S\,15.5 and
\S\,15.13~[27]: the operator is obliged not to cause harmful interference to
licensed radio services and to apply good engineering practice, and must
correct interference once it occurs, but no field strength or noise voltage
is prescribed. PLC systems themselves operate unlicensed under \S\,15.113 in
the band 9\,kHz to 490\,kHz on the utility's own lines; their frequency
assignments are coordinated between utilities through the PLC database of
the Utilities Technology Council (UTC)~[28]. In Europe the position is the
same in substance, with CISPR TR~18~[31] as the technical basis and no binding
limit for high-voltage installations.

The limits in a project specification are therefore design values set by the
owner. The methodology for converter-generated noise is that of CIGRE
TB~139~[29] and its successor TB~553~[30], which treat PLC and RI noise as
part of the AC filter design; the methodology for radio interference from
high-voltage installations, including the measurement procedure and the
derivation of limits from the protection of broadcast reception, is that of
CISPR TR~18-2~[31]. A clause written on this basis has five elements: the
frequency range, the limit, the measuring point and distance, the weather
condition and statistic, and the detector and bandwidth. A clause that omits
one of them cannot be verified against the supplier's design or on site.

Three conditions of the western United States change the values relative to
a European or eastern-US template. The corona onset of high-voltage hardware
falls with air density: at 1500\,m the relative air density is 0.84, and a
bundle or a fitting that meets a corona criterion at sea level does not meet
it at that elevation, so every corona-related value is calculated and
verified at the site or route elevation. The climate is dry, and for DC
corona the fair-weather emission of the positive pole is the governing case,
so the limits are stated for fair weather rather than for the foul-weather
case that AC line specifications emphasise. The population along a route is
sparse, with a small number of communities, while the interconnection areas
carry a dense set of third-party transmission lines on which carrier
protection remains in service; the PLC requirement is therefore written to
protect those third-party carriers, and the RI requirement is referred to the
station boundary and the edge of the right-of-way rather than to the nearest
dwelling.
\FloatBarrier
